\section{Discussion}
\label{sec:discussion}

In these experiments, CIF is best supported as a top-$k$ feature-selection
method for a downstream model. The benchmark asks whether a ranking places
predictive features early enough for fixed downstream learners, and CIF does so
at the level of random forests in classification. It is weaker at identifying
the informative features themselves: its synthetic recovery is in the lower
half, and its rankings are weakest at small $k$. Classification gives the
clearest support; the smaller regression panel is positive but less precise,
and its ordering is descriptive. The fixed-node
results provide the local guarantee; fitted tree and forest rankings are
assessed empirically.

\paragraph{Guidance.}
A CIF ranking should be judged through the downstream model that will use the
selected features. In practice, choose $k$ on held-out data over a small set of
values that includes the intended feature budget and, when feasible, the full
feature set as a sanity check. When prediction accuracy is the main target,
tree-ensemble baselines such as random forests, ExtraTrees, and boosted trees
should remain in the comparison set. When the target is variable screening or
a ranked feature list for later modeling, CIF should also be compared with
ctree, cforest, and forest-based screening methods such as Boruta, permutation
or conditional permutation importance, and RF-RFE
\citep{GenuerPoggiTuleauMalot2010VariableSelectionUsingRandomForests}.

The ablations locate most of the runtime in the threshold search. The
ablations show that threshold scanning and the 256-bin histogram account for
the largest measured CIF runtime changes, with small downstream score changes
in these experiments. On wide data the Stage~A test also matters, because the feature
correction scales every feature's budget. Adaptive stopping is inert at the
minimum budget and pays off only above it. The two corrections are part of the
testing rule and are not runtime shortcuts: removing them changes the trees
that are grown. For ranking quality, forest size matters more than the
split-count and feature-muting variants tested here, so the selected CIF
configuration is the relevant ranking reference when ranking quality is the
priority.

The max-type rule is an opt-in alternative to the per-threshold rule. It
reduces the reverse cardinality bias and costs less once a node has more than
about 50 candidates. It changes which features the trees rank but moves the
downstream scores only slightly (\cref{tab:maxt-benchmark-rankings}). The
benchmark rankings use the default per-threshold rule.

In sparse high-$p$ settings, feature use should be monitored directly.
The simulations show that CIF feature sampling can leave informative features
out of many forest split decisions, especially in sparse classification designs.
Useful diagnostics include top-$k$ stability across seeds, the share of splits
using the final top-$k$ features or any prespecified positive control features,
the number of distinct low-ranked or known null features used in splits, and
sensitivity to larger sampled feature sets or to a small CIF-all run with fewer
seeds or trees.

\paragraph{Limitations.}
The fixed-node results cover one node whose sample, tested features, and
permutation budget are fixed independently of the node responses, for
subsampled or unsampled fits. The benchmarked trees and forests add
adaptive stopping, feature and threshold ordering, feature muting, bounded
threshold sets, response-dependent choice of nodes, the Stage~A winner as the
feature tested at Stage~B, bootstrap samples, and forest aggregation, so their
behavior is assessed through the ranking benchmark. In a fitted tree the stored
Stage~B score is a screening statistic, not a post-selection $P$-value.

The benchmark coverage is uneven across analyses. Some comparisons use broader
panels where both methods are available, while complete-case rank and interval
summaries use smaller panels that include every method. These summaries are
therefore descriptive rather than confirmatory.

Configurations were selected on the reporting datasets with unequal grid sizes
(four for CIF, five for XGBoost, one for RF, ExtraTrees, CatBoost, and RF-RFE).
Selection optimism is therefore larger for the tuned families, and the order of
CIF and RF within the bunched group is not resolved. LODO measures the observed
sensitivity rather than bounding it.

The cardinality measurements use the linear selectors, while three of the four
selected configurations use RDC, so the reverse cardinality bias is established
for the linear selectors only. CPI ranks last in both tasks after a correction
to its correlation threshold, and its position should be read as provisional.
The three downstream learners are untuned, and the comparator set omits
Lasso and mutual-information filters.

The sampled-forest analysis isolates a sparse high-$p$ setting in which
forest feature sampling can limit how often splits use informative features:
larger sampled feature sets increase informative-feature use, but they also
change runtime and the mix of noise features used by splits. Track how often
splits use informative features alongside downstream performance and runtime
before changing the number of features sampled at each split.
